%% ============================================================
%%  SOLUCIÓN — Ejercicio 1.3: Listas
%%  Clase 1 — Después de diapositiva "Secciones y jerarquía"
%%  Compilar: F5 (pdflatex)
%% ============================================================

\documentclass[12pt, a4paper]{article}
\usepackage[utf8]{inputenc}
\usepackage[T1]{fontenc}
\usepackage[spanish]{babel}

\title{Fuentes de datos y metodología}
\author{Tu Nombre}
\date{\today}

\begin{document}
\maketitle

\section{Fuentes de datos}

Las principales fuentes de datos son:
\begin{itemize}
  \item Encuesta Permanente de Hogares (EPH)
  \item Censo Nacional de Población
  \item Registros administrativos de la DNIT
\end{itemize}

\section{Etapas de la investigación}

Las etapas de la investigación son:
\begin{enumerate}
  \item Recopilación y limpieza de datos
  \item Análisis descriptivo
  \item Estimación econométrica
  \item Interpretación y conclusiones
\end{enumerate}

\section{Ventajas de este enfoque}

Entre las ventajas de usar múltiples fuentes están:
\begin{itemize}
  \item Validación cruzada de resultados
  \item Mayor cobertura geográfica
  \item Representatividad nacional
  \item Información complementaria en diferentes dimensiones
\end{itemize}

\end{document}
